\subsubsection{Boundary and Bulk Correlation Scaling}
We next resolve the correlations along the wall without first averaging over
the transverse direction.  For a single trajectory $\xi$, define
\begin{align}
 C_{G,\xi}(x,r_y)
 &=\frac{1}{2N_y}\sum_{y,\mu,\nu}
 \left|C^{(\xi)}_{(x,y,\mu),(x,y+r_y,\nu)}\right|^2,
 \nonumber\\
 C_{G,\xi}^{\mathrm{av}}(r_y)
 &=\frac{1}{N_x}\sum_{x=0}^{N_x-1}C_{G,\xi}(x,r_y).
 \label{eq:EndpointSquaredCorrelator}
\end{align}
Here $C^{(\xi)}_{ij}=\operatorname{Tr}(\hat\rho_\xi\hat c_i^\dagger\hat c_j)$,
the $y$ coordinate is periodic, and both orbital indices are summed.  The
overline below denotes an arithmetic average over the independent
trajectories, after forming the squared correlator.  Thus we do not square
the averaged two-point function or average its logarithm.

This positive squared correlator has a direct density interpretation.
Write $\hat n_i=\hat c_i^\dagger\hat c_i$ and
$\langle\delta\hat n_i\delta\hat n_j\rangle_\xi
=\langle\hat n_i\hat n_j\rangle_\xi
-\langle\hat n_i\rangle_\xi\langle\hat n_j\rangle_\xi$.
Wick's theorem for a charge-conserving Gaussian trajectory gives
\begin{equation}
 \langle\delta\hat n_i\delta\hat n_j\rangle_\xi
 =\delta_{ij}C^{(\xi)}_{ii}-|C^{(\xi)}_{ij}|^2.
 \label{eq:EndpointDensityWickIdentity}
\end{equation}
For distinct modes, the connected density correlator is therefore the
negative of the squared two-point function.  With the total cell density
$\hat n_{x,y}=\sum_\mu\hat n_{x,y,\mu}$, the precise normalization of
our observable is
\begin{align}
 C_{G,\xi}(x,r_y)=-\frac{1}{2N_y}\sum_y
 \langle\delta\hat n_{x,y}\delta\hat n_{x,y+r_y}\rangle_\xi,
 \nonumber\\[-2pt]
 r_y\not\equiv0\pmod{N_y}.
 \label{eq:EndpointCellDensityCorrelator}
\end{align}
The factor $1/2$ is our per-orbital normalization.  At coincident cells the
contact term in Eq.~\eqref{eq:EndpointDensityWickIdentity} must be retained.
The connected subtraction is made inside each trajectory before the
ensemble average.  Forming a connected correlator of the mixed ensemble
instead adds the classical trajectory covariance of the local mean
densities, which is not included in the plotted quantity.

For a free chiral fermion, the fermion field has scaling dimension $1/2$,
so its two-point function decays as $d_{N_y}^{-1}$ and the squared
correlator has the critical prediction
\begin{equation}
 \overline{C_G^{\mathrm{av}}(r_y)}\sim A\,d_{N_y}(r_y)^{-\beta},
 \qquad \beta_{\mathrm{crit}}=2,
 \label{eq:EndpointCorrelatorCriticalPrediction}
\end{equation}
provided the long-distance signal is boundary dominated.  The chord is
$d_{N_y}(r_y)=(N_y/\pi)\sin(\pi r_y/N_y)$.  Figure~\ref{fig:HardWallCorrelatorSummary}
shows the distinction directly: at $\alpha_1=1$ the long-distance signal
survives the average over $x$, whereas at $\alpha_1=3$ it drops rapidly.
At $N_y=60$, the exact wall columns carry the largest long-distance
correlations.  The neighboring columns and the slab center have much
smaller amplitudes.  For this dataset we use the saved array coordinates
$x=0,\ldots,19$, with $x_L=5$ and $x_R=15$.

For the size collapse, we remove each circumference's amplitude using its
measured antipodal value.  With $r_\star=N_y/2$, set
\begin{align}
 X_N(r_y)&=\log\frac{d_{N_y}(r_y)}{d_{N_y}(r_\star)},
 \nonumber\\
 Y_N(r_y)&=\log\frac{\overline{C_G^{\mathrm{av}}(r_y)}}
                         {\overline{C_G^{\mathrm{av}}(r_\star)}}.
 \label{eq:EndpointCorrelatorCollapse}
\end{align}
Both coordinates vanish at the measured anchor; no fitted vertical shift
is introduced.  We fit $Y_N=-\beta X_N$ on
$8\leq r_y\leq N_y/2$, assigning equal total weight to each
circumference.  Explicitly, if $\mathcal W_N$ contains the $n_N$ selected
separations,
\begin{equation}
 \beta=-\frac{\sum_N n_N^{-1}\sum_{r_y\in\mathcal W_N}X_NY_N}
                  {\sum_N n_N^{-1}\sum_{r_y\in\mathcal W_N}X_N^2}.
 \label{eq:EndpointCorrelatorJointFit}
\end{equation}
The six-size fit gives $\beta=2.1903$, without constraining the power to
two.  We use $r_y\geq8$ to separate this fit from the short-distance
transient, while retaining the full independent interval through $N_y/2$.
The selected windows contain 5, 7, 9, 13, 18, and 23 separation values for
the six circumferences; $N_y=60$ does not supply 60 independent distances.
Moving the lower endpoint to five gives $\beta=2.2016$.
The earlier $2,\ldots,\lfloor N_y/4\rfloor$ window gave $2.2291$;
extending that same lower cutoff to $N_y/2$ gave $2.2282$.
Thus the lower cutoff, rather than the upper one, accounts for most of
the shift.  The antipodal normalization is the same in all these checks.
The collapse supports an
algebraic boundary contribution, but these finite-size effective exponents
remain above the ideal prediction.  This varies the circumference at
fixed $N_x=20$, not both dimensions of the system.  The fit is to logarithms
of ensemble-mean curves, not an average of sample-wise exponents.

\begin{figure}[!htbp]
 \centering
 \includegraphics[width=\columnwidth]{hard_wall_correlator_summary_3x1.pdf}
 \caption{\textbf{Endpoint correlation scaling.}
 Hard walls, $N_x=20$, $\alpha_2=30$, $n_{\rm shell}=1$, $t=2N_y$;
 $S=100$ independent pure half-filled initial states per ensemble,
 Born-conditioned exterior preparation, raster-$y$ perfect correction,
 and complex128 arithmetic.
 (a) Full-$x$ average at $N_y=60$, comparing $\alpha_1=1$ and $3$.
 (b) $\alpha_1=1$, $N_y=60$: walls $x=5,15$, adjacent cells $x=6,14$,
 and slab center $x=10$.
 (c) Antipodally normalized collapse of the $\alpha_1=1$ means for
 $N_y=24,28,32,40,50,60$.  The black dashed line is the equal-size-weighted
 fit $Y_N=-\beta X_N$ on $8\leq r_y\leq N_y/2$, giving $\beta=2.1903$.
 Light gray covers the union of the size-specific fit ranges; darker gray
 marks their common interval.  Each size retains its color and open symbol
 at every separation, including points outside its fit window.
 The dashed continuation outside the fit ranges is an extrapolation.
 All means precede logarithms; no uncertainty bands are shown.
 Only $1\leq r_y\leq N_y/2$ is used.  The $10^{-8}$ cutoff in (a,b)
 is for display, not a certified numerical floor; it excludes no fitted
 points in (c).  Colored lines in (a,b) guide the eye.}
 \label{fig:HardWallCorrelatorSummary}
\end{figure}

The boundary and bulk contributions need not be disentangled by subtracting
entropies.  For this particular observable, the spatial bookkeeping is
exactly additive.  Partition the columns into
$B=\{5,6,14,15\}$ and its complement $U$.  Let $\overline C_B$ and
$\overline C_U$ be the correlators averaged within their respective
sets and over trajectories.  Then
\begin{equation}
 \overline{C_G^{\mathrm{av}}}
 =\frac{|B|}{N_x}\overline C_B
 +\frac{|U|}{N_x}\overline C_U.
 \label{eq:EndpointCorrelatorSpatialSum}
\end{equation}
No independence assumption is needed: Eq.~\eqref{eq:EndpointSquaredCorrelator}
only samples pairs at the same $x$.  Correlations between different columns
would belong to a different observable.  The complement $U$ includes the
slab interior and exterior, and is not assumed to be a pure bulk sector.
For $N_y=60$, the four selected boundary columns contribute $99.35\%$ of
the full-$x$ average at $r_y=5$ and $99.95\%$ at $r_y=30$.

If the regional curves can be approximated by algebraic and exponential
pieces, this sum has the form
\begin{equation}
 \overline{C_G^{\mathrm{av}}(r_y)}
 \simeq A_{\mathrm e}d_{N_y}(r_y)^{-\beta}
 +A_{\mathrm b}\left[e^{-r_y/\xi_C}+e^{-(N_y-r_y)/\xi_C}\right].
 \label{eq:EndpointCorrelatorMixedDecay}
\end{equation}
The coefficients include the spatial weights.  This is a phenomenological
description of the squared correlators, not an exact decomposition of the
single-particle amplitudes: splitting $C=C_{\mathrm e}+C_{\mathrm b}$
inside one matrix element would generate an interference term in $|C|^2$.
At $x=10$, the observed short-distance curve is well described by
$\log\overline{C_G(10,r_y)}=b-r_y/\xi_C$ on $r_y=2,\ldots,6$, with
$\xi_C=0.6552$ unit cells.  This is the decay length of the squared
correlator over that window, not a claim of purely exponential behavior
at every separation.  A small algebraic tail can coexist in the interior
at finite slab width.  Once the exponential contribution has decayed,
the positive boundary contribution controls the full-$x$ average even
though only a few columns carry most of its weight.
